\documentclass[11pt, a4paper]{../problems_template}

\begin{document}

\begin{titlepage}
  \begin{center}
	\Huge{Турнир городов}
  \end{center}
\end{titlepage}

\chapter{1980-1981. Номер 2}

\begin{exercise}
На бесконечной плоскости играют двое: один передвигает одну фишку-волка, другой - одну из $K$ фишек-овец. После хода волка ходит какая-нибудь из овец, затем, после следующего хода волка, опять какая-нибудь из овец и так далее. И волк, и овцы передвигаются за один ход в любую сторону не более чем на один метр. Верно ли, что для любого числа овец, участвующих в игре, существует такая первоначальная позиция, что волк не поймает ни одной овцы?
\end{exercise}


\chapter{1981-1982. Номер 3}

\begin{exercise}
Рассматривается последовательность
$$
1, \frac{1}{2}, \frac{1}{3}, \frac{1}{4}, \frac{1}{5}, \frac{1}{6}, \frac{1}{7}, \dots
$$
Существует ли арифметическая прогрессия
\begin{itemize}
\item длины 5
\item сколь угодно большой длины,составленная из членов этой последовательности?
\end{itemize}
\end{exercise}
 
\begin{Solution}
Существует для обоих пунктов. Возьмем дробь вида $\frac{1}{n!}$. Любая кратная ей дробь, имеющая вид
$\frac{k}{n!}$ будет членом последовательности если $k \leqslant n$. Таким образом получим арифметическую прогрессию длины $k$ с шагом $\frac{1}{n!}$.
\end{Solution}

\begin{exercise}
Квадрат разбит на $K^{2}$ равных квадратиков. Про некоторую ломаную известно, что она проходит через центры всех квадратиков (ломаная может пересекать сама себя). Каково минимальное число звеньев у этой ломаной?
\end{exercise}

\chapter{1982-1983. Номер 4}

\begin{exercise}
Числа от 1 до 1000 расставлены по окружности. Докажите, что их можно соединить 500 непересекающимися отрезками, разность чисел на концах которых (по модулю) не более 749.
\end{exercise}

\begin{Solution}
Распределим числа в два множества - от 251 до 750 в первом и все оставшиеся во втором. Пусть одно множество будет покрашено в зеленый цвет, а другое в синий. Заметим что мощности этих множеств равны и что разность любых двух чисел разного цвета удовлетворяет условию задачи. Очевидно, что найдутся два соседних числа разного цвета. Соединим их прямой. Выкинем эту пару из окружности. Действуя по аналогии соединим все числа непересекающимися отрезками. 
\end{Solution}

\chapter{1983-1984. Номер 5}

\begin{exercise}
Из листа клетчатой бумаги размером $29 \times 29$ клеточек вырезали 99 квадратиков $2 \times 2$ (режут по линиям). Докажите, что из оставшейся части листа можно вырезать еще хотя бы один такой же квадратик.
\end{exercise}

\begin{exercise}
Дана бесконечная клетчатая бумага со стороной клетки, равной единице. Расстоянием между двумя клетками называется длина кратчайшего пути ладьи от одной клетки до другой (считается путь центра ладьи). В какое наименьшее число красок нужно раскрасить доску (каждая клетка закрашивается одной краской), чтобы две клетки, находящиеся на расстоянии 6, были всегда окрашены разными красками?
\end{exercise}


\chapter{1984-1985. Номер 6}

\begin{exercise}
Имеется 68 монет, причем известно, что любые две монеты различаются по весу. За 100 взвешиваний на \newline двухчашечных весах без гирь найдите самую тяжелую и самую легкую монеты.
\end{exercise}

\begin{Solution}
Взвесим 34 пары монет и получим две группы, которые назовем "Легкая" и "Тяжелая". Очевидно, что самая легкая монета находится в "Легкой" группе. Ее можно найти за 33 взвешивания. Аналогично за 33 взвешивания можно найти самую тяжелую монету в "Тяжелой" группе. Всего $34 + 33 + 33 = 100$ взвешиваний.
\end{Solution}

\begin{exercise}
В таблицу $10 \times 10$ нужно записать в каком-то порядке цифры $0, 1, 2, 3, \dots, 9$ так, что каждая цифра встречалась бы 10 раз.
\begin{enumerate}
\item Можно ли это сделать так, чтобы в каждой строке и в каждом столбце встречалось не более четырех различных цифр?
\item Докажите, что найдется строка или столбец, в которой (в котором) встречается не меньше четырех различных чисел.
\end{enumerate}
\end{exercise}

\begin{Solution}
Будем доказывать от противного. Предположим, что в строке или столбце не может быть больше 3 различных чисел. Обозначим количество различных чисел в строке $i$ как $r_{i}$, а количество различных чисел в столбце $i$ - $c_{i}$. Оценим сумму
$$
\sum_{i = 1}^{10} r_{i} + \sum_{i = 1}^{10} c_{i} \leqslant 3 \cdot 10 + 3 \cdot 10 = 60
$$
Теперь посчитаем эту сумму вторым способом. Если число встречается в строке, то оно добавляет 1 к $r_{i}$, а если в столбце то 1 к $c_{i}$. Обозначим количество строк в которых встречается произвольное число $i$ как $R_{i}$, а количество столбцов для него же как $C_{i}$. Несложно видеть, что $R_{i} \cdot C_{i} \geqslant 12$. Значит $R_{i} + C_{i} \geqslant R_{i} + \frac{12}{R_{i}}$. Тогда легко найти минимум $R_{i} + C_{i}$, который равен 7. Значит
$$
\sum_{i = 1}^{10} R_{i} + C_{i} \geqslant 70
$$
Но выше мы показали, что он не превышает 60. Противоречие.
\end{Solution}

\chapter{1985-1986. Номер 7}

\begin{exercise}
Последовательность чисел $x_{1}, x_{2}, \dots$ такова, что $x_{1} = \frac{1}{2}$ и для всякого натурального $k$ выполнено равенство $x_{k+1} = x_{k}^{2} + x_{k}$. Найдите целую часть суммы
$$
\frac{1}{x_{1}+1} + \frac{1}{x_{2}+1} + \dots + \frac{1}{x_{100}+1}
$$
\end{exercise}

\begin{Solution}
Сначала заметим, что
$$
\frac{1}{x_{k+1}} = \frac{1}{x_{k}(x_{k} + 1)} \implies \frac{1}{x_{k+1}} = \frac{1}{x_{k}} - \frac{1}{x_{k} + 1} \implies  \frac{1}{x_{k} + 1} = \frac{1}{x_{k}} - \frac{1}{x_{k + 1}}
$$
Откуда следует, что
$$
\frac{1}{x_{1}+1} + \frac{1}{x_{2}+1} + \dots + \frac{1}{x_{100}+1} = \frac{1}{x_{1}} - \frac{1}{x_{2}} + \frac{1}{x_{2}} - \frac{1}{x_{3}} + \dots + \frac{1}{x_{100}} - \frac{1}{x_{101}} = \frac{1}{x_{1}} - \frac{1}{x_{101}} = 2 - \frac{1}{x_{101}}
$$
Несложно заметить, что $x_{3} > 1$ и что последовательность возрастает. Значит $\frac{1}{x_{101}} < 1$. То есть, целая часть суммы равна 1.
\end{Solution}

\chapter{1986-1987. Номер 8}

\begin{exercise}
Берутся всевозможные непустые подмножества из множества чисел $1, 2, 3, \dots, N$. Для каждого подмножества берется величина, обратная к произведению всех его чисел. Найдите сумму всех таких обратных величин.
\end{exercise}

\begin{Solution}
Несложно заметить, что требуемая величина представляется в виде
$$
\left(1 + \frac{1}{1}\right)\left(1 + \frac{1}{2}\right)\left(1 + \frac{1}{3}\right)\dots \left(1 + \frac{1}{N}\right) - 1 = N
$$
\end{Solution}

\chapter{1987-1988. Номер 9}

\begin{exercise}
Докажите, что существует бесконечно много пар натуральных чисел $a$ и $b$ таких, что $a^{2} + 1$ делится на $b$, а $b^{2} + 1$ делится на $a$.
\end{exercise}

\begin{Solution}
Вспомним тождество Кассини для чисел Фибоначчи.
$$
(-1)^{n} = F_{n+1}F_{n-1} - F_{n}^{2}
$$
Тогда подставляя $F_{n} = F_{n+1} - F_{n-1}$ получим
$$
(-1)^{n} = F_{n+1}F_{n-1} - (F_{n+1} - F_{n-1})^{2} = 3F_{n+1}F_{n-1}-F_{n+1}^{2} - F_{n-1}^{2}
$$
Отсюда, для четного $n$ получим, что подходит пара $(F_{n+1}, F_{n-1})$. 
\end{Solution}

\chapter{1988-1989. Номер 10}

\begin{exercise}
Числа $1, 2, 3, \dots, N$ записываются в строчку в таком порядке, что если где-то (не на первом месте) записано число $i$, тогде-то слева от него встретится хотя бы одно из чисел $i+1$ и $i-1$. Сколькими способами это можно сделать?
\end{exercise}

\chapter{1989-1990. Номер 11}

\begin{exercise}
Рассмотрим все возможные наборы чисел из множества $\{1, 2, 3, \dots, N\}$, не содержащие двух соседних чисел. Докажите, что сумма квадратов произведений чисел в этих наборах равна $(N+1)!-1$.
\end{exercise}

\begin{Solution}
Будем доказывать по индукции. База очевидна. Составим рекуррентное соотношение для искомой величины $S_{n}$. Она состоит из наборов которые включают $n$, самого $n$ и наборов, в которые $n$ не входит. Если $n$ входит в набор, то он учитывается в общем выражении как $n^{2}S_{n-2}$, а если нет, то используется $S_{n-1}$. В итоге
$$
S_{n} = n^{2}S_{n-2}+S_{n-1}+n^{2} = ((n-1)! - 1)n^{2} + (n! - 1) + n^{2} = n!(n + 1) - 1 = (n + 1)! - 1
$$  
\end{Solution}

\chapter{1995-1996. Номер 17}

\begin{exercise}
\begin{enumerate}
\item Восемь учеников решали восемь задач. Оказалось, что каждую задачу решили 5 учеников. Докажите, что найдутся такие два ученика, что каждую задачу решил хотя бы один из них.
\item Если каждую задачу решило 4 ученика, то может оказаться, что таких двоих не найдется. Докажите это.
\end{enumerate}
\end{exercise}

\begin{Solution}
Рассмотрим двудольный граф из учеников и задач, где наличие ребра между учеником и задачей означает, что ученик не решил задачу. Тогда с одной стороны существует $8 \binom{8 - 5}{2} = 24$ уголка в множестве задач. Но в множестве учеников $\binom{8}{2} = 28$ пар. Значит существует пара учеников не создающая ни одного уголка. А значит они вместе решили все задачи. 
\end{Solution}

\chapter{1996-1997. Номер 18}

\begin{exercise}
\begin{enumerate}
\item Квадрат разрезан на равные прямоугольные треугольники с катетами 3 и 4 каждый. Докажите, что число треугольников четно.
\item Прямоугольник разрезан на равные прямоугольные треугольники с катетами 1 и 2 каждый. Докажите, что число треугольников четно.
\end{enumerate}
\end{exercise}

\chapter{1997-1998. Номер 19}

\begin{exercise}
Шахматный король обошел всю доску $8 \times 8$, побывав на каждой клетке по одному разу, вернувшись последним ходом в исходную клетку. Докажите, что он сделал четное число диагональных ходов.
\end{exercise}

\begin{Solution}
При ходе по диагонали цвет поля на котором стоит король не меняется. Так как король вернулся обратно, то для каждого хода c черной клетки на белую должен найтись ход с белой клетки на черную. Значит ходов меняющих цвет клетки четное число. Всего король сделал 64 хода, а значит диагональных ходов тоже четное число. 
\end{Solution}

\begin{exercise}
\begin{itemize}
\item На доске выписаны числа 1, 2, 4, 8, 16, 32, 64, 128. Разрешается стереть любые два числа и вместо них выписать их разность - неотрицательное число. После семи таких операций на доске будет только одно число. Может ли оно равняться 97?
\item На доске выписаны числа $1, 2^{1}, 2^{2}, 2^{3}, \dots, 2^{10}$. Разрешается стереть любые два числа и вместо них выписать их разность - неотрицательное число. После нескольких таких операций на доске будет только одно число. Чему оно может быть равно?
\end{itemize}
\end{exercise}

\begin{Solution}
Ответ на первый пункт - да. 
$$
16 - 8 = 8, \quad 8 - 4 = 4, \quad 4 - 2 = 2, \quad 2 - 1 = 1, \quad 
64 - 32 = 32, \quad 32 - 1 = 31, \quad 128 - 31 = 97
$$
Заметим, что четность суммы всех чисел на доске не меняется. Значит последнее оставшееся число не может быть четным. Докажем по индукции, что можно получить любое нечетное число от 1 до $2^{10} - 1$. Если на доске только 1 и 2, то можно получить только 1. Теперь положим, что на доске числа от 1 до $2^{n+1}$ и из чисел от 1 до $2^{n}$ можно получить любое нечетное от 1 до $2^{n} - 1$. Тогда это можно сделать и с $2^{n+1}$ выполнив в качестве первой операции $2^{n+1} - 2^{n}$. Получить нечетные от $2^{n} + 1$ до $2^{n+1} - 1$ можно операцией $2^{n+1} - A$, где $A$ - любое нечетное, которое можно получить из чисел от 1 до $2^{n}$.
\end{Solution}

\begin{exercise}
Назовем лабиринтом шахматную доску $8 \times 8$, где между некоторыми полями вставлены перегородки. Если ладья может обойти все поля, не перепрыгивая через перегородки, то лабиринт называется хорошим, иначе - плохим. Каких лабиринтов больше - хороших или плохих?
\end{exercise}

\chapter{1998-1999. Номер 20}

\begin{exercise}
В круге провели $2n$ радиусов. Они разделили круг на $2n$ равных секторов: $n$ синих и $n$ красных, чередующихся в произвольном порядке. В синие сектора, начиная с некоторого, записывают против хода часовой стрелки числа от 1 до $n$. В красные сектора, начиная с некоторого, записывают те же числа, но по ходу часовой стрелки. Докажите, что найдется полукруг, в котором записаны все числа от 1 до $n$.
\end{exercise}

\begin{Solution}

\end{Solution}

\chapter{1999-2000. Номер 21}

\begin{exercise}
Несколько последовательных натуральных чисел выписали в строку в таком порядке, что сумма любой тройки подряд идущих чисел делится нацело на самое левое число этой тройки. Какое максимальное количество чисел могло быть выписано, если последнее число строки - нечетное?
\end{exercise}

\begin{Solution}
Заметим, что если в строке стоит два четных числа подряд, то следующие за ними числа тоже должны быть четными. Иначе бы возникла ситуация, когда подряд следуют четное, четное и нечетное числа. И это бы противоречило условию. Значит после каждого четного числа стоит два нечетных. Но тогда должно быть не более двух четных чисел. Иначе бы разность между количеством четных и нечетных чисел была бы не менее двух, что противоречит условию. Значит чисел не более 5. Примером может служить 2, 1, 3, 4, 5.
\end{Solution}

\chapter{2000-2001. Номер 22}

\begin{exercise}
В весеннем туре турнира городов 2000 года старшеклассникам страны $N$ было предложено 6 задач. Каждую задачу решили ровно 1000 школьников, но никакие два школьника не решили вместе все 6 задач. Каково наименьшее возможное число старшеклассников страны $N$, принявших участие в весеннем туре?
\end{exercise}

\chapter{2007-2008. Номер 29}

\begin{exercise}
\begin{enumerate}
\item Фокусник с завязанными глазами выдает зрителю 5 карточек с номерами от 1 до 5. Зритель прячет две карточки, а остальные отдает ассистенту фокусника. Ассистент показывает зрителю две из них, и зритель называет номера этих карточек фокуснику(в том порядке, в каком захочет). После этого фокусник угадывает номера карточек, спрятанных у зрителя. Как фокуснику и ассистенту договориться, чтобы фокус всегда удавался?
\item  То же самое, но карточек 29, с номерами от 1 до 29.
\end{enumerate}
\end{exercise}

\begin{exercise}
Одиннадцати мудрецам завязывают глаза и надевают каждому на голову колпак одного из 1000 цветов. После этого им глаза развязывают, и каждый видит все колпаки, кроме своего. Затем одновременно каждый показывает остальным одну из двух карточек - белую или чёрную. После этого все должны одновременно назвать цвет своих колпаков. Удастся ли это? Мудрецы могут заранее договориться о своих действиях (до того, как им завязали глаза); мудрецам известно, каких 1000 цветов могут быть колпаки. 
\end{exercise}


\chapter{2008-2009. Номер 30}

\begin{exercise}
В ячейку памяти компьютера записали число 6. Далее компьютер делает миллион шагов. На шаге номер $n$ он увеличивает число в ячейке на наибольший общий делитель этого числа и $n$. Докажите, что на любом шаге компьютер увеличивает число в ячейке либо на 1, либо на простое число.
\end{exercise}

\begin{Solution}

\end{Solution}

\chapter{2012-2013. Номер 34}

\begin{exercise}
Числа $1, 2, \dots, 100$ стоят по кругу в некотором порядке. Может ли случиться, что у любых двух соседних чисел модуль разности не меньше 30, но не больше 50?
\end{exercise}

\begin{Solution}
Разделим числа на две равные по размеру группы. В первую группу поместим числа от $n$ до $n + 49$, где $n$ выберем позже. Во вторую группу поместим все остальные числа. Выберем $n$ так, чтобы разность между любыми двумя числами во второй группе была либо меньше 30, либо больше 50. Условия для этого - $100 - (n + 50) < 30$ и $n - 2 < 30$. Значит подходит $20 < n < 32$. Выберем, например, $n = 22$. Итак, в первой группе числа от 22 до 71, а во второй остальные. Числа второй группы не могут соседствовать друг с другом, а значит они должны чередоваться с числами первой группы. Но с числом 22 может стоять только одно число второй группы - это число 72. Получили противоречие.
\end{Solution}

\chapter{2015-2016. Номер 37}

\begin{exercise}
Из целых чисел от 1 до 100 удалили $k$ чисел. Обязательно ли среди оставшихся чисел можно выбрать $k$ различных чисел с суммой 100, если
\begin{itemize}
\item $k = 9$
\item $k = 8$
\end{itemize}
\end{exercise}

\begin{Solution}
Пусть удалены числа от 1 до 9 включительно. Тогда минимальная сумма любых 9 из оставшихся чисел будет не меньше 
$$
\sum_{i = 0}^{8} (10 + i) = 126
$$
Поэтому для $k = 9$ ответ нет. Во втором случае рассмотрим пары чисел дающих в сумме 25. 1 и 24, 2 и 23, $\dots$, 12 и 13. Таких пар всего 12. Значит в любом случае после удаления останется не менее 4 пар. Сумма чисел в которых будет 100.
\end{Solution}

\chapter{2016-2017. Номер 38}

\begin{exercise}
Десяти ребятам положили в тарелки по 100 макаронин. Есть ребята не хотели и стали играть. Одним действием кто-то из детей перекладывает из своей тарелки по одной макаронине всем другим детям. После какого \\ наименьшего количества действий у всех в тарелках может оказаться разное количество макаронин?
\end{exercise}

\begin{exercise}
В каждой клетке доски $8 \times 8$ написали по одному натуральному числу. Оказалось, что при любом разрезании доски на доминошки суммы чисел во всех доминошках будут разные. Может ли оказаться, что наибольшее записанное на доске число не больше 32? Доминошкой называется прямоугольник, состоящий из двух клеток.
\end{exercise}

\begin{Solution}
Раскрасим доску шахматной раскраской. В черные клетки запишем числа $1, 2, 3, \dots, 32$, а в белые 1. Таким образом, при любом разрезании суммы чисел в доминошках будут $2, 3, 4, \dots, 33$.
\end{Solution}

\chapter{2017-2018. Номер 39}

\begin{exercise}
В каждой вершине выпуклого многогранника сходятся три грани. Каждая грань покрашена в красный, жёлтый или синий цвет. Докажите, что число вершин, в которых сходятся грани трёх разных цветов, чётно.
\end{exercise}

\begin{Solution}
Если вершина трехцветная, то пойдем по желто-синему ребру. Вершина на другом конце ребра либо трехцветная, либо имеет еще одно желто-синее ребро кроме того, по которому пришли в нее. Продолжая путь мы в итоге упремся в трехцветную вершину, которая очевидно не равна стартовой. Значит все трехцветные вершины разбиваются на пары. Значит их четное количество.
\end{Solution}

\chapter{2019-2020. Номер 41}

\begin{exercise}
Для каких $N$ можно расставить в клетках квадрата $N \times N$ действительные числа так, чтобы среди всевозможных сумм чисел на парах соседних по стороне клеток встречались все целые числа от $1$ до $2(N-1)N$ включительно (ровно по одному разу)?
\end{exercise}

\begin{Solution}
Существует для всех $N > 1$. Рассмотрим, например, квадрат для $N = 6$:
\begin{align*}
\begin{pmatrix}
1  &  1 & 3  &  3 & 5  & 5  \\
11 & 13 & 13 & 15 & 15 & 17 \\
23 & 23 & 25 & 25 & 27 & 27 \\
33 & 35 & 35 & 37 & 37 & 39 \\
45 & 45 & 47 & 47 & 49 & 49 \\
55 & 57 & 57 & 59 & 59 & 61
\end{pmatrix}
\end{align*}
Каждый элемент выбирается так, чтобы суммы соседних элементов в строке или столбце образовывали последовательность $2, 4, 6, 8, 10, 12, 14, \dots, 120$. Тогда, разделив каждый элемент квадрата на два, мы получим искомое. Для нечетных $N$ квадрат строится аналогично.
\end{Solution}


\chapter{2020-2021. Номер 42}

\begin{exercise}
Существуют ли 100 таких натуральных чисел, среди которых нет одинаковых, что куб одного из них равен сумме кубов остальных?
\end{exercise}

\chapter{2021-2022. Номер 43}

\begin{exercise}
На доске написано число 7. Петя и Вася по очереди приписывают к текущему числу по одной цифре, начинает Петя. Цифру можно приписать в начало числа (кроме нуля), в его конец или между любыми двумя цифрами. Побеждает тот, после чьего хода число на доске станет точным квадратом. Может ли кто-нибудь гарантированно победить, как бы ни играл соперник?
\end{exercise}

\begin{exercise}
На столе в ряд лежат 20 плюшек с сахаром и 20 с корицей в произвольном порядке. Малыш и Карлсон берут их по очереди, начинает Малыш. За ход можно взять одну плюшку с любого края. Малыш хочет, чтобы ему в итоге досталось по десять плюшек каждого вида, а Карлсон пытается ему помешать. При любом ли начальном расположении плюшек Малыш может достичь своей цели, как бы ни действовал Карлсон?
\end{exercise}

\begin{exercise}
Султан собрал 300 мудрецов и предложил им испытание. Он сообщил им список из 25
цветов и сказал, что на испытании каждому мудрецу наденут на голову колпак одного
из этих цветов, причём если для каждого цвета написать количество надетых колпаков этого цвета, все числа будут различны. Каждый мудрец увидит, какой колпак на ком надет, но свой колпак не увидит. Затем одновременно (по сигналу) каждый должен будет назвать предполагаемый цвет своего колпака. Могут ли мудрецы заранее договориться действовать так, чтобы гарантированно хотя бы 150 из них назвали цвет верно?
\end{exercise}

\chapter{2022-2023. Номер 44}

\begin{exercise}
Натуральные числа от 1 до 100 раскрашены в три цвета: 50 чисел - в красный, 25
чисел - в жёлтый и 25 чисел - в зелёный. Известно, что все красные и жёлтые числа можно разбить на 25 троек так, чтобы в каждой тройке было два красных числа и одно жёлтое, которое больше одного красного и меньше другого. Аналогичное утверждение верно для красных и зелёных чисел. Обязательно ли все 100 чисел можно разбить на 25 четвёрок, в каждой из которых два красных числа, одно жёлтое и одно зелёное, при этом жёлтое и зелёное числа лежат между красными?
\end{exercise}

\begin{Solution}
Обозначим $i$-тое по порядку желтое число - $y_{i}$, зеленое - $g_{i}$, а красные разделим на две группы. Первые по порядку 25 - $r_{i}$, а оставшиеся - $R_{i}$. Очевидно $y_{i} > r_{i}$, в ином случае для какого-то желтого не нашлось бы меньшего красного. Аналогично - $r_{i} < g_{i}$. Те же рассуждения доказывают, что $y_{i} < R_{i}$ и $g_{i} < R_{i}$. Значит мы можем выбрать в качестве $i$-той четверки $r_{i}$, $g_{i}$, $y_{i}$, $R_{i}$.
\end{Solution}

\chapter{2023-2024. Номер 45}

\begin{exercise}
Имеется кучка из 100 камней. Играют двое. Первый берёт 1 камень, потом второй берёт 1 или 2 камня, потом первый берёт 1, 2 или 3 камня, затем второй 1, 2, 3 или 4 камня и так далее. Выигрывает взявший последний камень. Кто может обеспечить себе победу, как бы ни играл его соперник?
\end{exercise}

\begin{Solution}
Выигрывает первый. Представим 100 как сумму $1 + 3 + 5 + 7 + \dots + 19$. Обозначим последовательность взятых камней $1, s_{1}, f_{1}, s_{2}, f_{2}, \dots$ и $f_{i} \leqslant 2i + 1$, а $s_{i} \leqslant 2i$. Тогда, если $f_{i} = 2i + 1 - s_{i}$, то перед последними двумя ходами останется 19 камней, и второй игрок не сможет взять их все, а первый сможет взять оставшиеся. 
\end{Solution}

\chapter{2024-2025. Номер 46}

\begin{exercise}
Учитель назвал две различные ненулевые цифры. Коля хочет составить делящееся на 7 семизначное число, в десятичной записи которого нет других цифр, кроме этих двух. Всегда ли Коля может это сделать, какие бы две цифры ни назвал учитель?
\end{exercise}

\begin{exercise}
Назовём ходы коня, при которых он смещается на две клетки по горизонтали и на одну по вертикали, \\ горизонтальными, а остальные - вертикальными. Требуется поставить коня на одну из клеток доски $46 \times 46$, после чего чередовать им горизонтальные и вертикальные ходы. Докажите, что если запрещено посещать клетки более одного раза, то будет сделано не более 2024 ходов.
\end{exercise}

\begin{exercise}
Даны две строго возрастающие последовательности положительных чисел, в которых каждый член, начиная с третьего, равен сумме двух предыдущих. Известно, что каждая последовательность содержит хотя бы одно число, которого нет в другой последовательности. Какое наибольшее количество общих чисел может быть у этих последовательностей?
\end{exercise}


\end{document} 
 
 
